\section{Messbarkeit des Nutzens} \label{sec:messbarkeit}

\subsection{Metriken zur Leistungsbewertung} \label{sec:metriken}

Für die Bewertung KI-gestützter Bildanalyse sind standardisierte Klassifikationsmetriken
etabliert, die im Sinne der SMART-Kriterien eingesetzt werden. Die vier zentralen
Messgrößen sind:

\begin{itemize}
    \item \textbf{Sensitivität (Recall):} Anteil der korrekt erkannten positiven Fälle.
          In der medizinischen Diagnostik ist eine hohe Sensitivität besonders relevant,
          da übersehene Befunde unmittelbare klinische Konsequenzen haben können.
    \item \textbf{Spezifität:} Anteil der korrekt erkannten negativen Fälle. Falsch
          positive Befunde erzeugen unnötige Folgeuntersuchungen.
    \item \textbf{F1-Score:} Harmonisches Mittel aus Precision und Recall, besonders
          relevant bei unausgewogener Klassenverteilung -- etwa bei seltenen
          Krankheitsbildern.
    \item \textbf{AUC-ROC:} Fläche unter der Receiver-Operating-Characteristic-Kurve;
          klassenunabhängiges Gesamtmaß der Klassifikationsleistung, unabhängig vom
          gewählten Entscheidungsschwellenwert.
\end{itemize}

Diese Metriken gelten für industrielle und medizinische KI-Systeme gleichermaßen und
unterstreichen die methodische Konvergenz beider Domänen.

\subsection{KI im Vergleich zum Menschen} \label{sec:ki_vs_mensch}

Der Standardbenchmark in der medizinischen KI-Forschung ist der direkte Vergleich
zwischen KI-System und menschlichen Fachärzten auf einem gemeinsamen, verblindeten
Testdatensatz. Tabelle~\ref{tab:ki_vergleich} fasst aktuelle Ergebnisse aus der Literatur
zusammen:

\begin{table}[htbp]
    \centering
    \caption{Leistungsvergleich KI vs.\ Radiologen in ausgewählten Studien}
    \label{tab:ki_vergleich}
    \begin{tabularx}{\textwidth}{|X|c|c|l|}
        \hline
        \textbf{Anwendung} & \textbf{Sensitivität} & \textbf{Spezifität} & \textbf{Einordnung} \\
        \hline
        Pneumonie (Röntgen)        & 88\,\%            & 90\,\%               & vergleichbar mit Radiologen \\
        Lungenkarzinom (Röntgen)   & 85\,\%            & 88\,\%               & vergleichbar mit Radiologen \\
        Krebsdiagnose (KI-assist.) & 79\,\%            & 87\,\%               & besser als ohne Unterstützung \\
        Brustkrebsscreening        & \(\times\)1{,}12  & \(\approx\)1{,}00    & +12\,\% Sens., --44\,\% Aufwand \\
        \hline
    \end{tabularx}
    \source{\cite{metaAnalysisChest2025, breastScreening2025}}
\end{table}

Die Ergebnisse zeigen konsistent: Der Einsatz von KI als ergänzender zweiter Leser
erhöht die Sensitivität bei minimalem Spezifitätsverlust. Der Nutzenzuwachs entsteht
nicht durch Substitution des Radiologen, sondern durch Kombination: Das System reduziert
den Befundungsaufwand und markiert auffällige Regionen; die abschließende klinische
Entscheidung verbleibt beim Arzt \cite{aiSecondReader2024}.
